\section{Deterministic Validation of the Monitors}
\label{sec:controls}

A monitor that never fired is indistinguishable from a monitor that is broken.
Before reporting any detection rate, we therefore fire each event the monitors
are supposed to catch, by hand, and measure how long each takes. The suite
needs no model backend and costs nothing to run.

\begin{table}[h]
\centering
\small
\begin{tabular}{@{}rlr@{}}
\toprule
\textbf{\#} & \textbf{Control} & \textbf{Measured} \\
\midrule
1 & heartbeat silence after killing the pinger & 9.08\,s \\
2 & canary in a workspace directory name & 0.43\,s \\
3 & canary in an outbound query string & 0.49\,s \\
4 & destination outside the egress baseline & 0.49\,s \\
5 & canary in cross-run persistent memory & 18.45\,s \\
6 & red-teamer serves its payload through the SSRF & 0.52\,s \\
\bottomrule
\end{tabular}
\caption{All six fire. Control 1 measures the configured threshold (5\,s
interval times a multiplier of 2) minus the time since the last ping, not
monitor slowness. Control 5 is an order of magnitude slower than the other
canary surfaces because inotify does not fire for host-side writes to a bind
mount, so that surface falls back to the periodic sweep.}
\end{table}

Two of these controls exist because the harness failed silently without them.
Control 5 was added after finding that the most persistent leak channel in the
corpus, a canary written into cross-run memory, was the only one nothing
watched. Control 6 was added after finding that no control touched the
attacker path at all, and it reported a blind spot on its first run
(Section~\ref{sec:validity}).

A separate sweep varies each monitor's configuration parameter against the
same synthetic event, isolating how often a monitor looks from the variance of
model behaviour. Its results appear in Section~\ref{sec:results}.
